\documentclass[10.5pt,a4paper]{article}
\usepackage[top=16mm,bottom=16mm,left=22mm,right=22mm]{geometry}
\usepackage{xeCJK}
\setCJKmainfont{Microsoft YaHei}
\setmainfont{Segoe UI}
\usepackage{xcolor}
\usepackage{titlesec}
\usepackage{enumitem}
\usepackage{hyperref}

\definecolor{primary}{HTML}{1a56db}

\hypersetup{colorlinks=true,urlcolor=primary,linkcolor=primary}

\titleformat{\section}{\large\bfseries}{}{0em}{}
\titleformat{\subsection}{\normalsize\bfseries}{}{0em}{}
\titlespacing{\section}{0pt}{7pt}{3pt}
\titlespacing{\subsection}{0pt}{5pt}{2pt}

\setlength{\parindent}{2em}
\setlength{\parskip}{2pt}
\pagestyle{empty}

\begin{document}

\begin{center}
{\LARGE\bfseries MoonEvidence 项目申报书}
\end{center}

\section{基本信息}

\begin{itemize}[leftmargin=16pt,itemsep=1pt,parsep=0pt]
  \item 项目名称：MoonEvidence：可信证据包验证平台
  \item 参赛者：陈俊文
  \item 联系方式：187****1181
  \item GitHub 仓库链接：\url{https://github.com/starlittle/MoonEvidence}
  \item Gitlink 仓库链接：\url{https://gitlink.org.cn/starlittle/MoonEvidence}
  \item 项目方向：MoonBit 数据完整性验证平台 / 工程基础设施
  \item 是否为移植项目：否（原创项目）
  \item 许可证：Apache-2.0
\end{itemize}

\section{项目简介}

MoonEvidence 是一个纯 MoonBit 实现的可信证据包验证库和 CLI 工具。给定一组文件及其元数据、SHA-256/SHA-512 摘要、Merkle 树证明、Ed25519 数字签名和版本链记录，MoonEvidence 能验证这些内容是否完整且未被篡改，并输出结构化的诊断报告。项目还提供 Trust Workbench 交互式工作台，支持可视化 Merkle 树和篡改实验室（Tamper Lab）。

项目定位为链无关验证核心——不是区块链应用或智能合约框架，而是一个可被任何需要证明"这组文件没被改过"的场景复用的通用验证引擎。核心逻辑零 IO 依赖，同一套代码可运行于 CLI、CI 和浏览器三种环境。

\section{核心功能范围}

\begin{itemize}[leftmargin=16pt,itemsep=1pt,parsep=0pt]
  \item 提供 RFC 8785 Canonical JSON 序列化，确保 JSON 输出字节级稳定，可安全用于摘要计算；完整实现 ECMAScript 最短数字表示和 code-unit 级别键排序；
  \item 提供纯 MoonBit SHA-256 和 SHA-512 双算法摘要实现，通过 NIST 标准测试向量验证，js 后端 SHA-256 吞吐量约 58 MiB/s；
  \item 提供 RFC 6962 风格 Merkle 树构建与证明验证，带域分离前缀，支持双算法，与独立 Node 参考实现交叉校验；
  \item 提供证据清单（manifest）模型和验证，拒绝路径遍历攻击，支持文件级摘要比对和 Merkle 根校验；
  \item 提供线性版本链验证，确保版本记录的时间线性和哈希链完整性；
  \item 提供 7 步验证流水线，从解析、规范化、摘要比对到 Merkle 根重算的完整管线，输出带冻结错误码（E1xxx--E5xxx、W1xxx）的结构化 findings 报告；
  \item 提供 Ed25519 数字签名，从 GF(2\textsuperscript{255}-19) 有限域到 RFC 8032 签名/验签全自实现，通过 RFC 8032 §7.1 已知答案测试，支持审计日志条目签名；
  \item 提供内容寻址存储（SHA-256 对象存储），支持去重和证据包重建；
  \item 提供追加式哈希链审计日志，每条记录链接前一条哈希，支持 Ed25519 签名验证，篡改可检测；
  \item 提供 CLI 工具 \texttt{verify [-{}-json]}、\texttt{explain} 和 \texttt{create}，冻结退出码（0 通过 / 1 失败 / 2 错误），支持机器可读 JSON 报告和人类可读诊断文本；
  \item 提供浏览器端验证适配器，同一核心编译为 ESM 模块，纯客户端验证无需上传服务器，配套 Trust Workbench 6 视图交互式工作台（验证 / 创建 / 证明 / 审计 / 签名 + 篡改实验室），支持 Merkle 树可视化和实时篡改对比；
  \item 三后端支持：native / wasm-gc / js 全部通过构建和测试，GitHub Actions CI 三后端矩阵构建；
  \item 提供 265 个单元测试、41 个黑盒 CLI 测试、属性测试（规范化幂等性、Merkle 证明健全性）、变异验证测试和 1000 节点链压测，配套 criterion 风格性能基准（全流水线约 26 $\mu$s/file）；代码规模 9884 行（实现 4866 + 测试 5018），88 个有效 commit，12 个包；
  \item 提供 README 示例，覆盖 CLI 验证、JSON 报告、explain 输出和浏览器端使用。
\end{itemize}

\section{原创说明}

本项目为原创实现，非移植项目。设计参考了以下公开标准：

\begin{itemize}[leftmargin=16pt,itemsep=1pt,parsep=0pt]
  \item RFC 8785（JSON Canonicalization Scheme）
  \item RFC 6962（Certificate Transparency, Merkle Tree 部分）
  \item RFC 8032（Ed25519 数字签名）
  \item NIST FIPS 180-4（SHA-256 / SHA-512）
\end{itemize}

项目在 MoonBit 生态中填补数据完整性验证这一空白，实现了 Mooncakes 生态中首个 RFC 8785 Canonical JSON 序列化，可作为其他项目的可信数据基础层发布至 Mooncakes 包管理器。

\end{document}
